% Level 3 answer key, Week 8 -- Data Detectives
\answersweek{8}{Data Detectives}{Answers for workbook chapter 8 and homework 8.}

\section{Exercise 8.1: Survey cleaner}
\pyfile[firstline=12,lastline=45]{../excercises/week08/ex-survey-cleaner/solution.py}
\begin{answer}[Answers]
Mess in \texttt{favourite\_games.csv}: spaces in \texttt{" block craft "}, \texttt{"Kart Dash "} and \texttt{" 5 "}; capitals in \texttt{BLOCK CRAFT}, \texttt{block craft}, \texttt{kart dash}; empty row 5 and a row of only commas; missing hours in row 7.
The cleaner throws away \textbf{2 rows}: response 5 and the all-commas row, because they have no favourite game (19 raw rows $\rightarrow$ 17 clean). Row 7 is kept: its game is useful even though the hours are missing (\py|None|).
\end{answer}

\section{Exercise 8.2: Survey stats}
\pyfile[firstline=40,lastline=65]{../excercises/week08/ex-survey-stats/solution.py}
\begin{answer}[Answers]
Most popular game: \textbf{Block Craft} (5 votes). Average hours per week: \textbf{3.1} (50 hours $\div$ 16 answers; the missing answer is skipped). Average stars: 4.4.
Adventure (5.0 stars from 2 answers): \textbf{do not trust it much} -- the sample is tiny; one more answer could change it a lot. Any answer mentioning \enquote{only 2 people} is correct.
\end{answer}

\section{Exercise 8.3: Bar chart}
\pyfile[firstline=31,lastline=70]{../excercises/week08/ex-bar-chart/solution.py}
\begin{answer}[What to look for]
Six bars, tallest first: Block Craft 5, Kart Dash 4, Puzzle Drop 3, Sky Islands 2, Pixel Pets 2, Space Chess 1. Every bar has its number on top and its name underneath; heights are in proportion (the 5-vote bar is 5 times the 1-vote bar). Ask: \enquote{Could someone who cannot see colours read this chart?} -- the answer should be yes.
\end{answer}

\section{Reflection prompts}
\begin{answer}[Possible answers]
The chart tells us which games this group likes most. It \emph{cannot} tell us what all children like (only 17 made-up answers), why they like them, or whether people answered honestly. On the survey question: good answers mention privacy and choosing not to share a full name.
\end{answer}

\section{Homework 8}
\pyfile[firstline=15,lastline=43]{../excercises/week08/hw-survey-privacy/solution.py}
\begin{answer}[Core]
\texttt{club} is risky because a small club plus an answer (\enquote{the chess-club person who reads 35 minutes}) can point to one person even without a name. The first anonymous row is \texttt{\{'id': 'R1', 'favourite\_snack': 'apple slices', 'minutes\_reading': '20'\}}.
\end{answer}
\begin{answer}[Challenge]
The anonymous file has three columns (\texttt{id}, \texttt{favourite\_snack}, \texttt{minutes\_reading}) and 6 rows. Most popular snack: a \textbf{tie} between apple slices and popcorn (2 each). Praise students who notice the tie rather than naming just one.
\end{answer}
\begin{answer}[Creative]
Look for: a consent sentence, no names/ages/addresses/schools, a question with a number answer, and questions that are kind (nothing embarrassing).
\end{answer}
